\documentclass[11pt,a4paper]{article}
\usepackage[margin=1.1in]{geometry}
\usepackage{amsmath,amssymb}
\usepackage[colorlinks=true,allcolors=blue]{hyperref}
\usepackage{parskip}
\usepackage{xr}
\externaldocument{build/main}

% Generated by make_numbers.py -- do not edit by hand.
\newcommand{\AblBaseline}{-0.0277}
\newcommand{\AblBestCell}{uniform / derived / all / tuned}
\newcommand{\AblBestRtwo}{0.7807}
\newcommand{\AblFull}{0.7622}
\newcommand{\AblNcells}{32}
\newcommand{\AblNorders}{24}
\newcommand{\CfgManualFull}{0.6925}
\newcommand{\CfgManualSub}{0.6901}
\newcommand{\CfgManualTopFull}{0.7438}
\newcommand{\CfgManualTopSub}{0.7155}
\newcommand{\CfgUniformFull}{0.7453}
\newcommand{\CfgUniformSub}{0.7423}
\newcommand{\CfgUniformTopFull}{0.7664}
\newcommand{\CfgUnscaledFull}{-0.0277}
\newcommand{\CfgUnscaledSub}{-0.0538}
\newcommand{\CvCorrhi}{0.719}
\newcommand{\CvCorrlo}{0.667}
\newcommand{\CvMean}{0.6931}
\newcommand{\CvSd}{0.0368}
\newcommand{\CvVOnehi}{0.765}
\newcommand{\CvVOnelo}{0.621}
\newcommand{\DSAmesBase}{0.6819}
\newcommand{\DSAmesFull}{0.8358}
\newcommand{\DSAmesFullUni}{0.8381}
\newcommand{\DSAmesN}{1022}
\newcommand{\DSAmesP}{36}
\newcommand{\DSAmesScal}{-0.1558}
\newcommand{\DSAmesSelect}{+0.1606}
\newcommand{\DSAmesTune}{+0.1491}
\newcommand{\DSAmesUniScal}{+0.1211}
\newcommand{\DSAmesUniSelect}{+0.0068}
\newcommand{\DSAmesUniTune}{+0.0283}
\newcommand{\DSCaliforniaBase}{0.5013}
\newcommand{\DSCaliforniaFull}{0.7563}
\newcommand{\DSCaliforniaFullUni}{0.7768}
\newcommand{\DSCaliforniaN}{10000}
\newcommand{\DSCaliforniaP}{18}
\newcommand{\DSCaliforniaScal}{+0.1616}
\newcommand{\DSCaliforniaSelect}{+0.0186}
\newcommand{\DSCaliforniaTune}{+0.0749}
\newcommand{\DSCaliforniaUniScal}{+0.1869}
\newcommand{\DSCaliforniaUniSelect}{+0.0157}
\newcommand{\DSCaliforniaUniTune}{+0.0730}
\newcommand{\DSKcHouseBase}{0.3888}
\newcommand{\DSKcHouseFull}{0.8092}
\newcommand{\DSKcHouseFullUni}{0.7806}
\newcommand{\DSKcHouseN}{10000}
\newcommand{\DSKcHouseP}{20}
\newcommand{\DSKcHouseScal}{+0.0805}
\newcommand{\DSKcHouseSelect}{+0.1197}
\newcommand{\DSKcHouseTune}{+0.2202}
\newcommand{\DSKcHouseUniScal}{+0.2710}
\newcommand{\DSKcHouseUniSelect}{+0.0053}
\newcommand{\DSKcHouseUniTune}{+0.1155}
\newcommand{\FlipA}{SVR uniform / raw8 / default / n=3000}
\newcommand{\FlipB}{SVR manual / top12 / tuned / full}
\newcommand{\FlipDiff}{-0.0271}
\newcommand{\FlipP}{0.166}
\newcommand{\FlipPuncorr}{0.007}
\newcommand{\GbmFull}{0.8384}
\newcommand{\GbmSmall}{0.8055}
\newcommand{\GridFits}{192}
\newcommand{\GridRtwo}{0.7806}
\newcommand{\GridSecs}{422}
\newcommand{\GridSecsMin}{7.0}
\newcommand{\KbelowPlateau}{0.5745}
\newcommand{\KbestAt}{10}
\newcommand{\KbestK}{10}
\newcommand{\KernLinear}{0.6469}
\newcommand{\KernPolyThree}{0.6650}
\newcommand{\KernPolyTwo}{0.6613}
\newcommand{\KernRbf}{0.7702}
\newcommand{\KernSigmoid}{0.3207}
\newcommand{\KmaxAt}{18}
\newcommand{\KminAt}{4}
\newcommand{\KnnFull}{0.7599}
\newcommand{\KnnSmall}{0.7179}
\newcommand{\KplateauFrom}{10}
\newcommand{\KplateauRange}{0.0045}
\newcommand{\KrangeHi}{0.7683}
\newcommand{\KrangeLo}{0.5745}
\newcommand{\MaxDeriv}{+0.5428}
\newcommand{\MaxScal}{+0.7202}
\newcommand{\MaxSelect}{+0.5937}
\newcommand{\MaxTune}{+0.5312}
\newcommand{\MinDeriv}{-0.0881}
\newcommand{\MinScal}{-0.0239}
\newcommand{\MinSelect}{-0.0029}
\newcommand{\MinTune}{+0.0152}
\newcommand{\NDatasets}{3}
\newcommand{\NFlipped}{2}
\newcommand{\NModels}{10}
\newcommand{\NOrderHoldsAuto}{1}
\newcommand{\NOrderHoldsUni}{3}
\newcommand{\NPairs}{21}
\newcommand{\NestedRandom}{0.7532}
\newcommand{\NestedRandomMax}{0.767}
\newcommand{\NestedRandomMin}{0.744}
\newcommand{\NestedSpatial}{-0.5344}
\newcommand{\NestedSpatialMax}{0.660}
\newcommand{\NestedSpatialMin}{-4.091}
\newcommand{\NystroemM}{2000}
\newcommand{\NystroemRtwo}{0.7673}
\newcommand{\NystroemSecs}{18}
\newcommand{\PairManualDiff}{+0.0242}
\newcommand{\PairManualP}{4.6\times 10^{-6}}
\newcommand{\PairManualPuncorr}{3.4\times 10^{-9}}
\newcommand{\PairRfDiff}{-0.0446}
\newcommand{\PairRfP}{9.3\times 10^{-5}}
\newcommand{\PairRfPuncorr}{9.4\times 10^{-8}}
\newcommand{\PairXgbDiff}{-0.0728}
\newcommand{\PairXgbP}{1.9\times 10^{-7}}
\newcommand{\PairXgbPuncorr}{1.2\times 10^{-10}}
\newcommand{\PreethiReportedMSE}{1.12}
\newcommand{\PreethiReportedRtwo}{0.15}
\newcommand{\RandSixtyRtwo}{0.7834}
\newcommand{\RandSixtySecs}{203}
\newcommand{\RandTwentyRtwo}{0.7702}
\newcommand{\RandTwentySecs}{100}
\newcommand{\RandTwoHundredRtwo}{0.7796}
\newcommand{\RandTwoHundredSecs}{1377}
\newcommand{\RandomCVmean}{0.7496}
\newcommand{\RangeDeriv}{0.6309}
\newcommand{\RangeScal}{0.7441}
\newcommand{\RangeSelect}{0.5966}
\newcommand{\RangeTune}{0.5160}
\newcommand{\RepScaledDefault}{0.737}
\newcommand{\RepScaledTuned}{0.766}
\newcommand{\RepScaledTunedMSE}{0.309}
\newcommand{\RepSeeds}{5}
\newcommand{\RepUnscaledDefault}{-0.020}
\newcommand{\RepUnscaledDefaultMSE}{1.343}
\newcommand{\RepUnscaledTuned}{0.514}
\newcommand{\RepUnscaledTunedMSE}{0.640}
\newcommand{\RfFull}{0.8171}
\newcommand{\RfSmall}{0.7610}
\newcommand{\RidgeFull}{0.6504}
\newcommand{\RidgeSmall}{0.6497}
\newcommand{\SearchBest}{random 60}
\newcommand{\SearchBestRtwo}{0.7834}
\newcommand{\SearchBestSecs}{203}
\newcommand{\SelectionOptimism}{0.0000}
\newcommand{\ShapDeriv}{0.1505}
\newcommand{\ShapScal}{0.2723}
\newcommand{\ShapSelect}{0.1789}
\newcommand{\ShapTune}{0.1882}
\newcommand{\SizeFull}{0.7813}
\newcommand{\SizeFullN}{14448}
\newcommand{\SizeFullRmse}{0.5382}
\newcommand{\SizeFullSecs}{8.3}
\newcommand{\SizeSmall}{0.6775}
\newcommand{\SizeSmallN}{1000}
\newcommand{\SizeThreeK}{0.7418}
\newcommand{\SizeThreeKrmse}{0.5848}
\newcommand{\SpatialCVmean}{0.2316}
\newcommand{\SpatialCVsd}{0.7525}
\newcommand{\SpatialOptimism}{0.5180}
\newcommand{\StatRho}{0.429}
\newcommand{\StatSplits}{10}
\newcommand{\StatSvrCIhi}{0.783}
\newcommand{\StatSvrCIlo}{0.769}
\newcommand{\StatSvrMean}{0.7757}
\newcommand{\StatSvrNBhi}{0.792}
\newcommand{\StatSvrNBlo}{0.759}
\newcommand{\StatSvrSd}{0.0099}
\newcommand{\StatSvrVOnehi}{0.795}
\newcommand{\StatSvrVOnelo}{0.756}
\newcommand{\SvrRankFull}{4}
\newcommand{\SvrRankSmall}{4}
\newcommand{\SvrRbfFull}{0.7702}
\newcommand{\SvrRbfSmall}{0.7501}
\newcommand{\ThreshRange}{0.0004}
\newcommand{\WeightMax}{0.7674}
\newcommand{\WeightMin}{0.7128}
\newcommand{\WeightMinOverlap}{10}
\newcommand{\WeightN}{66}
\newcommand{\WeightRange}{0.0546}
\newcommand{\XgbFull}{0.8483}
\newcommand{\XgbSmall}{0.8057}


\setlength{\emergencystretch}{2em}
\pagestyle{empty}

\begin{document}

\noindent
Emmanuel Adutwum\\
Soka University of America\\
Aliso Viejo, California, United States\\
\texttt{eadutwum@soka.edu}

\vspace{14pt}
\noindent 3 September 2026

\vspace{14pt}
\noindent
Dr Santosh Reddy Addula, Editor\\
Mr Ramakant Mohite, Assistant Editor\\
\emph{Discover Artificial Intelligence}\\
Springer Nature

\vspace{14pt}
\noindent
\textbf{Re: Submission ID 6cf4e30e-1a5b-4ad2-99e5-2320f0737507 --- revised manuscript}

\vspace{10pt}
\noindent Dear Dr Addula and Mr Mohite,

Please find enclosed the revised version of my manuscript, now entitled ``A Controlled
Reassessment of Support Vector Regression on the California Housing Benchmark with Feature
Scaling, Feature Selection and Hyperparameter Tuning'', submitted in response to the decision
of major revision. A point-by-point response to every comment from the Editor and all three
reviewers accompanies this letter, together with an updated Contribution Information Sheet.

I want to be direct about what changed, because the revision is not cosmetic. Two of the
central claims of the previous version were wrong, and the reviewers found both.

First, the value of $R^2 \approx 0.60$ that the previous version attributed to the SVR of
Preethi et al.\ is not their SVR result at all --- it is the score of their linear, ridge and
polynomial-ridge models. Their SVR-RBF result is approximately \PreethiReportedRtwo{}.
Reviewer~1 and Reviewer~3 independently observed that an unscaled baseline of $-0.054$ could
not explain a reported $0.60$, and they were right: the premise was mistaken. Reproducing the
configuration they describe, over \RepSeeds{} seeds, now yields $R^2 = \RepUnscaledDefault$
with library defaults and $\RepUnscaledTuned$ when tuned --- values that bracket their actual
result --- while adding a single standardisation step to the identical pipeline yields
$\RepScaledTuned$. The causal story the paper tells is now supported by a direct reproduction
rather than by an assumed correspondence.

Second, the previous version's claim that feature scaling accounted for 95.7\% of the
improvement was an artefact of evaluating a single ablation ordering. Evaluating all
\AblNcells{} cells of the four-factor design shows that the apparent contribution of scaling
varies by \RangeScal{} in $R^2$ depending only on where it is placed in the sequence. The
order-independent Shapley value is $\ShapScal$. Scaling remains the largest single
contributor, but the previous percentage is withdrawn.

The revision also reports two results that contradict the earlier framing and that I have
chosen to report rather than omit. The feature-specific scaling strategy I previously
presented as the paper's preprocessing advance does \emph{not} outperform plain
standardisation on this dataset --- it is worse, by a difference that is significant under a
paired test. And the derived features contribute little once selection and tuning are in
place. In addition, spatially blocked cross-validation, which Reviewer~1 requested, reveals
that the optimism of random partitioning on these geographically structured data exceeds
every pipeline effect measured anywhere in the study. That finding qualifies the headline
numbers of this literature, including my own, and it is now stated as such.

Beyond these, the revision adds a direct reproduction of the prior configuration, an
equal-budget comparison in which all \NModels{} models receive identical data and identical
tuning, corrected confidence intervals with paired significance testing on identical splits
under the Nadeau--Bengio correction, feature selection repeated inside every fold together
with fully nested estimates, replication of the workflow on King County and Ames housing
data, and sensitivity analyses for every hand-chosen constant. The workflow itself is now
defined explicitly as a staged procedure with a diagram, the methods sections have been
expanded substantially, the claim of novelty has been withdrawn and replaced with a
methodological and diagnostic contribution, and the manuscript has been restructured and
copy-edited throughout.

On the reporting inconsistencies noted by Reviewer~1, I took a structural approach rather
than a proofreading one. Every experiment now writes a JSON record, and every table, figure
and in-text number in the manuscript is generated from those records by script. No numerical
value is typed by hand anywhere in the paper, so the text, the tables and the figures cannot
diverge.

Following the internal editorial feedback, the Declarations section now addresses ethical
approval, consent to participate and consent to publish under individual subheadings; the
data availability statement appears above the References with direct links to every dataset
and to the code repository; and a dual publication declaration records the arXiv preprint
(arXiv:2605.08660) and the relationship between it and the present manuscript. As Reviewer~2
requested and the Editor endorsed, the manuscript now includes a disclosure of my use of a
large language model as an assistive tool, together with a statement of what remains my own
responsibility.

I am grateful to the reviewers. The requirement to reproduce the prior configuration directly
is what uncovered my misreading of the reference value, and the requirement to evaluate the
full ablation design is what showed that my headline attribution depended on an arbitrary
ordering. The paper is more modest than it was, and considerably more solid.

The manuscript has not been published elsewhere and is not under consideration by any other
journal. I have no competing interests to declare. I would be glad to provide any further
information the Editor or reviewers may require.

\vspace{12pt}
\noindent Thank you for your consideration.

\vspace{18pt}
\noindent Sincerely,

\vspace{22pt}
\noindent Emmanuel Adutwum\\
Soka University of America

\end{document}
